% ch-a2-01.tex —— 第 2 篇（Xinwen Zhu）篇首 + §1 + §2.1–2.2 前半
% 底本 PDF 52–64 页 = 印刷页 39–51
\chapter{参数叠上的凝聚层}
\ysrc{39}
\begin{center}
{\small\itshape Xinwen Zhu}
\end{center}
\medskip
\noindent\textbf{摘要.} 我们在局部（$\ell\ne p$）与整体函数域两种情形中构造算术 Langlands 参数的叠。我们提出若干关于这些叠上某些假想凝聚层的猜想，并解释它们在局部与整体 Langlands 纲领中所起的作用。我们综述若干已知结果，作为这些猜想的证据。
\footnote{2020 年数学主题分类：主分类 11S37，次分类 11R39、11G18。}
\footnote{作者的工作部分受 NSF 资助（协议号 DMS-1902239），并获 Simons 奖学金资助。}

\medskip
\noindent\textbf{目录.}
\begin{list}{}{\leftmargin=2.2em\labelwidth=2em\labelsep=0.4em\itemindent=-2.2em}
\item 引言
\item 表示空间
\item 算术 Langlands 参数的叠
\item 参数叠上的凝聚层
\item 致谢
\item 参考文献
\end{list}

\section{引言}

近些年来，以下一点变得越来越清楚：在局部与整体算术 Langlands 参数的叠上，应当存在某些凝聚层 $\mathcal A$（确切地说，是凝聚层的复形）。人们期待这些层在很大程度上支配 Langlands 对应，并使人能够在算术 Langlands 纲领内部表述局部-整体相容性。这类对象的存在性，已经在族化局部 Langlands 对应的框架下，由 Emerton–Helm [25] 与 Helm [39] 的工作所提示。\footnote{在 $p$-进局部 Langlands 纲领中也出现类似的 $\mathcal A$，讨论见 [24]。} 这一想法最近由 Hellmann [38] 进一步发展。另一方面，在 V.~Lafforgue 与 Genestier–Lafforgue [36, 49] 的工作之后，这些想法变得更加具体，而几何 Langlands 纲领中的有力工具如今已可用于实现其中的（一部分）。事实上，人们期待非阿基米德局部域上的整个算术局部 Langlands 对应容许一个范畴化的化身（例如见 [32, 4.2] 中的若干提示）。这类凝聚层的存在性自然地嵌入这一范畴框架，而这正是本文所要解释的。
\ysrc{40}
与此相关的一个方向上，Fargues–Scholze [27] 关于局部 Langlands 对应几何化的工作也与这些想法密切一致，并且同样指向算术局部 Langlands 对应的一个范畴形式。从整体的视角看，$\mathcal A$ 的存在性在作者与 Xiao [77] 的合作工作中起着指导原则的作用，该工作通过志村簇的上同调给出 Jacquet–Langlands 对应的几何实现。在另一个方向上，这类凝聚层的一种粗糙形式出现在作者与 V.~Lafforgue [50] 的合作工作中，那里用它来描述 Shtuka 上同调的椭圆部分，所处框架是 Arthur–Kottwitz 猜想。

在本文中，我们提出若干与这些假想层 $\mathcal A$ 相关的精确猜想，并综述已知结果，其中包括在某些特殊（但特别重要）的情形下对 $\mathcal A$ 的显式猜想性描述，以及它们在局部-整体相容性中的作用。我们还提出局部算术 Langlands 对应的一个猜想性范畴形式，它为这类 $\mathcal A$ 的预期存在性提供概念上的依据。为了表述这些猜想，我们讨论局部 Langlands 参数（在 $\ell\ne p$ 的情形）与整体 Langlands 参数（在函数域情形）的模叠的构造及其若干性质。我们注意到，本文中的一些想法已在专家之间非正式地流传多年。\footnote{事实上，在本文第一版公开前后，出现了若干相关工作。例如见 [3, 5, 19, 27, 38]。} 作者的意图是把其中一些想法表述得更精确，并把它们写下来。

本文自然地分为两部分。第 2 节与第 3 节致力于表示模空间的一般研究以及 Langlands 参数模空间的构造。由于这些节中的结果是新的，我们对几乎所有的断言都给出详细证明。第 4 节专门表述我们的主要猜想。它包含一些新的结果（例如定理 4.45），对此我们同样给出详细证明。同时，这一节还综述那些为我们的猜想提供证据的已知或即将发表的结果，因此在若干处带有更强的综述性质。

\section{表示空间}

设 $M$ 是交换环 $\Lambda$ 上的仿射群概形，$\Gamma$ 是一个抽象群。众所周知，存在 $\Lambda$ 上的仿射概形 ${}^{cl}\mathcal R_{\Gamma,M}$，使得对每个 $\Lambda$-代数 $A$，集合 ${}^{cl}\mathcal R_{\Gamma,M}(A)$ 分类从 $\Gamma$ 到 $M(A)$ 的群同态。具体地说，先考虑 $\Lambda$ 上分类从 $\Gamma$ 到 $M(A)$ 的一切集合映射的函子。这个函子显然由 $M$ 在 $\Gamma$ 上的自乘积 $M^\Gamma$ 表示。再加上「这些集合映射是群同态」这一条件，就把 ${}^{cl}\mathcal R_{\Gamma,M}$ 定义为 $M^\Gamma$ 的闭子概形。

第一个问题众所周知：伽罗瓦群是投射有限群，人们必须考虑它们的连续表示，并附加某些额外性质。我们将在第 2.4 节处理这个问题。粗略地说，施加连续性条件之后，我们得到的 ind-概形，其在闭点处的完备化还原投射有限群表示的通常加框架形变空间。然而，这样的空间一般可能不具有好的整体几何（见例 2.50）。尽管如此，在第 3 节所考虑的情形中，这些空间以合理的方式把所有形变空间「粘」在一起。\ysrc{41}

第二个问题在于：定义 ${}^{cl}\mathcal R_{\Gamma,M}\subset M^\Gamma$ 的方程通常不构成「正则序列」，这可能使 ${}^{cl}\mathcal R_{\Gamma,M}$ 上出现非平凡的导出结构。在下文的若干处，我们需要记住这些空间可能的导出结构。因此，我们将在第 2.2 节回顾导出对象的构造。这一构造当然众所周知（例如见 [35, 69]），但我们的处理方式将受 [49] 启发，并在第 2.1 节先回顾幺半群的导出范畴。

\subsection{幺半群的导出范畴}

我们的目标是定义一个导出几何对象 $\mathcal R_{\Gamma,M}$，用以参数化从 $\Gamma$ 到 $M$ 的同态。为此，先从一个更一般的框架出发、考虑幺半群的同态是方便的。想法是从幺半群的范畴 $\mathbf{Mon}$ 走到它的导出范畴。由于 $\mathbf{Mon}$ 非阿贝尔，我们必须采用 Quillen 意义上的非阿贝尔导出范畴这一概念，它由 Lurie 用 $\infty$-范畴的语言发展起来 [54, 5.5.8]。我们先回顾一些一般理论，然后特殊化到与本文相关的例子。

下文中，我们把 $(\infty,1)$-范畴简称为 $\infty$-范畴，并按通常方式把普通范畴视为 $\infty$-范畴。以 $\mathbf{Spc}$ 记空间的 $\infty$-范畴，它把集合的范畴 $\mathbf{Sets}$ 作为全子范畴包含在内（把集合视为离散空间）。包含 $\mathbf{Sets}\to\mathbf{Spc}$ 有左伴随 $\pi_0:\mathbf{Spc}\to\mathbf{Sets}$，它保持有限积。若 $x,y$ 是 $\infty$-范畴 $\mathcal C$ 中的两个对象，我们记 $\Map_{\mathcal C}(x,y)\in\mathbf{Spc}$ 为从 $x$ 到 $y$ 的映射空间。（即使 $\mathcal C$ 是普通范畴我们也使用这个记号，此时该空间是离散的。）所有函子都在 $\infty$-范畴的意义下理解（因而是导出的）。以 $\mathbf{Fun}(\mathcal C,\mathcal D)$ 记两个 $\infty$-范畴 $\mathcal C$ 与 $\mathcal D$ 之间函子的 $\infty$-范畴。以 $\Delta$ 记（普通的）单纯范畴。关于 $\infty$-范畴的基础，我们引用 [54]。

我们认为，从 Clausen–Scholze 的观点出发是有启发性的。对容许余极限的普通范畴 $\mathcal C$，以 $\mathcal C^{cp}$ 记其紧投射对象构成的全子范畴，即那些 $x\in\mathcal C$，使得 $\Map_{\mathcal C}(x,-)$ 与滤过余极限及自反余等化子交换。这是一个容许有限余积的范畴，于是可以定义其非阿贝尔导出范畴 $\mathcal P_\Sigma(\mathcal C^{cp})$（[54, 5.5.8.8]），即 $\mathbf{Fun}((\mathcal C^{cp})^{\mathrm{op}},\mathbf{Spc})$ 中由保持有限积的函子构成的全子范畴。\footnote{我们默认 $\mathcal C^{cp}$ 是小范畴，我们所遇到的一切例子都如此。} 若 $\mathcal C$ 由 $\mathcal C^{cp}$ 在 1-范畴余极限下生成（非正式地说，这意味着 $\mathcal C$ 中的对象可由 $\mathcal C^{cp}$ 中的对象通过取「并」（滤过余极限）与「表示」（自反余等化子）得到），则 $\mathcal P_\Sigma(\mathcal C^{cp})$ 被 Clausen–Scholze 称为 $\mathcal C$ 的 anima 的 $\infty$-范畴，记作 $\mathbf{Ani}(\mathcal C)$。（一般性叙述见 [14, §5]。）我们有时也直接称它为 $\mathcal C$ 的导出范畴。$\mathbf{Ani}(\mathcal C)$ 中的对象可由 $\mathcal C^{cp}$ 在 $\infty$-范畴余极限下生成。注意，若 $\mathcal C$ 带有对称幺半结构，使得张量积分别对每个变量保持余极限，且该对称幺半结构限制为 $\mathcal C^{cp}$ 上的对称幺半结构，则 $\mathbf{Ani}(\mathcal C)$ 自然地成为对称幺半 $\infty$-范畴，且张量积分别对每个变量保持余极限（[55, 4.8.1.10]）。

存在全忠实嵌入 $\mathcal C\subset\mathbf{Ani}(\mathcal C)$：把 $\mathcal C$ 视为保持有限积、且分解为 $(\mathcal C^{cp})^{\mathrm{op}}\to\mathbf{Sets}\subset\mathbf{Spc}$ 的函子 $(\mathcal C^{cp})^{\mathrm{op}}\to\mathbf{Spc}$ 构成的范畴。它有余伴随 $\pi_0:\mathbf{Ani}(\mathcal C)\to\mathcal C$，由 $\pi_0:\mathbf{Spc}\to\mathbf{Sets}$ 诱导。\ysrc{42}

更一般地，对每个 $m\ge 0$，有 $m$-截断函子 $\tau_{\le m}:\mathbf{Ani}(\mathcal C)\to{}_{\le m}\mathbf{Ani}(\mathcal C)$，其中对 $\infty$-范畴 $\mathcal C$，${}_{\le m}\mathcal C$ 记 $\mathcal C$ 中 $m$-截断对象构成的全子范畴（[54, 5.5.6.1]），它是自然包含函子 ${}_{\le m}\mathbf{Ani}(\mathcal C)\subset\mathbf{Ani}(\mathcal C)$ 的左伴随（[54, 5.5.6.18]）。下面是几个基本例子。

\begin{exam}{2.1}\label{exam:a2-2-1}
(1) 若 $\mathcal C=\mathbf{Sets}$，带笛卡尔对称幺半结构（即张量积由积给出），则 $\mathcal C^{cp}$ 是有限集的范畴 $\mathbf{Sets}_f$，且 $\mathbf{Ani}:=\mathbf{Ani}(\mathbf{Sets})\cong\mathbf{Spc}$（[54, 5.5.8.24]），带笛卡尔对称幺半结构。由于这一自然的等价，下文中我们将不加区分地使用 $\mathbf{Ani}$ 与 $\mathbf{Spc}$。

(2) 设 $\Lambda$ 是交换环。若 $\mathcal C=\mathbf{Mod}_\Lambda^\heartsuit$ 是 $\Lambda$-模的阿贝尔范畴，带通常的张量积结构，则 $\mathcal C^{cp}$ 是有限投射 $\Lambda$-模的范畴，而 $\mathbf{Ani}(\mathbf{Mod}_\Lambda^\heartsuit)$ 等价于 $\Lambda$-模的连通复形（即上同调在正次数消没的复形\footnote{本文对 $\Lambda$-模的稳定 $\infty$-范畴 $\mathbf{Mod}_\Lambda$ 中的复形采用上同调约定。于是对 $N\in\mathbf{Mod}_\Lambda$，我们写 $H^iN=\pi_{-i}N$，并对满足 $H^i(N[j])=H^{i+j}N$ 的对象记作 $N[j]$。同调代数中通常的截断函子写作 $\tau^{\le n},\tau^{\ge n}:\mathbf{Mod}_\Lambda\to\mathbf{Mod}_\Lambda$，这与 [54, 5.5.6.18] 中的截断函子 $\tau_{\le m}$ 不同。不过，$\tau^{\ge -m}$ 在 $\mathbf{Mod}_\Lambda^{\le 0}$ 上的限制同构于 $\tau_{\le m}$。}）的导出范畴 $\mathbf{Mod}_\Lambda^{\le 0}:=D^{\le 0}(\mathbf{Mod}_\Lambda^\heartsuit)$，带通常的对称幺半结构（[54, 5.5.8.21] 与 [14, 5.1.6]）。
\end{exam}

我们需要的例子是幺半群的范畴 $\mathcal C=\mathbf{Mon}$。这个范畴容许所有小余极限，并由其紧投射对象在余极限下生成，这些紧投射对象正是有限自由生成的幺半群。对有限集 $I$，以 $\mathrm{FM}(I)$ 记由 $I$ 生成的自由幺半群。以 $\mathbf{FFM}$ 记由这些 $\mathrm{FM}(I)$ 张成的全子范畴。对幺半群 $\Gamma$，以 $\mathbf{FFM}/\Gamma$ 记相应的切片范畴：即对象是形如 $(\mathrm{FM}(I),u:\mathrm{FM}(I)\to\Gamma)$ 的对，从 $(\mathrm{FM}(I),u)$ 到 $(\mathrm{FM}(J),v)$ 的态射是幺半群同态 $f:\mathrm{FM}(I)\to\mathrm{FM}(J)$ 且满足 $u=vf$。我们注意到，范畴 $\mathbf{FFM}/\Gamma$ 不是滤过的，但是筛状的（这一概念见 [54, 5.5.8.1]），因为 $\mathbf{FFM}/\Gamma$ 中存在余积。在 $\mathbf{Mon}$ 中有典范同构
\begin{equation}\tag{2.1}
\varinjlim_{\mathbf{FFM}/\Gamma}\mathrm{FM}(I)\xrightarrow{\ \sim\ }\Gamma .
\end{equation}
这一同构也可以在 $\mathbf{Ani}(\mathbf{Mon})$ 中理解：经由全忠实嵌入 $\mathbf{Mon}\subset\mathbf{Ani}(\mathbf{Mon})$，有 $\mathbf{Ani}(\mathbf{Mon})=\mathcal P_\Sigma(\mathbf{FFM})$。

另一方面，对容许有限积的 $\infty$-范畴 $\mathcal C$，有 $\mathcal C$ 中幺半群对象的 $\infty$-范畴 $\mathbf{Mon}(\mathcal C)$，按定义它是 $\mathcal C$ 中单纯对象的范畴
\[
\mathcal C_\Delta:=\mathbf{Fun}(\Delta^{\mathrm{op}},\mathcal C)
\]
的全子范畴，由那些 $X_\bullet$ 组成：对每个 $[n]\in\Delta$，由 $[1]\cong\{i-1,i\}\subset\{0,1,\dots,n\}=[n]$ 诱导的映射
\[
X([n])\to X(\{0,1\})\times X(\{1,2\})\times\cdots\times X(\{n-1,n\})=X([1])^n
\]
是 $\mathcal C$ 中的同构（[55, 4.1.2.5]）。例如，若 $\mathcal C=\mathbf{Sets}$，则由通常的 Milnor 构造有 $\mathbf{Mon}\cong\mathbf{Mon}(\mathbf{Sets})$：对 $\Gamma\in\mathbf{Mon}$，$\mathbf{Mon}(\mathbf{Sets})$ 中相应的对象是「恰有一个对象、其自同态幺半群为 $\Gamma$」的范畴的脉（[55, 4.1.2.4]）。于是全忠实嵌入 $\mathbf{Sets}\subset\mathbf{Spc}$ 诱导全忠实嵌入 $\mathbf{Mon}\subset\mathbf{Mon}(\mathbf{Spc})$（两者都是 $\mathbf{Spc}_\Delta$ 的全子范畴）。\ysrc{43}

\begin{exam}{2.2}\label{exam:a2-2-2}
回忆存在从（小）$\infty$-范畴的 $\infty$-范畴到 $\mathbf{Spc}_\Delta$ 的全忠实嵌入，它把 $\mathcal C$ 送到单纯空间 $[n]\mapsto\mathbf{Fun}(\Delta^n,\mathcal C)^{\simeq}$，即 $\mathbf{Fun}(\Delta^n,\mathcal C)$ 中最大的 Kan 复形。这个函子的本质像由所谓的完全 Segal 空间组成。这样，每个只有一个对象的 $\infty$-范畴都给出 $\mathbf{Spc}$ 中的一个幺半群对象。由 [55, 4.7.1]，给定 $\infty$-范畴 $\mathcal C$ 中的对象 $x$，存在幺半群 $\End_{\mathcal C}(x)\in\mathbf{Mon}(\mathbf{Spc})$，其在 $[1]\in\Delta$ 处的值同构于 $\Map_{\mathcal C}(x,x)$，并且它在所有作用于 $x$ 的 $\mathbf{Mon}(\mathbf{Spc})$ 对象中具有泛性质。我们称它为 $x$ 的导出自同态幺半群。
\end{exam}

现在我们有两个可以看作 $\mathbf{Mon}$ 的导出版本的 $\infty$-范畴。幸运的是，它们典范等价。

\begin{lem}{2.3}\label{lem:a2-2-3}
存在典范等价 $\mathbf{Ani}(\mathbf{Mon})\cong\mathbf{Mon}(\mathbf{Spc})$。
\end{lem}

\begin{proof}
我们考虑更一般的情形。设 $\mathcal C$ 是如前所述的（普通）余完备对称幺半范畴（即 $\mathcal C$ 由 $\mathcal C^{cp}$ 在余极限下生成，且张量积分别对每个变量保持余极限）。那么谈论它在 $\mathcal C$ 与 $\mathbf{Ani}(\mathcal C)$ 中的结合（即 $E_1$-）代数对象的（$\infty$-）范畴 $\mathbf{Alg}(-)$ 是有意义的（[55, 2.1.3]）。利用 [55, 7.2.4.27] 与下面的引理 2.4，我们得到典范等价
\[
\mathbf{Ani}(\mathbf{Alg}(\mathcal C))\cong\mathbf{Alg}(\mathbf{Ani}(\mathcal C)).
\]
令 $\mathcal C=\mathbf{Sets}$，并在环绕的对称幺半结构是笛卡尔结构时把结合代数对象与幺半群对象等同起来（[55, 2.4.2, 4.1.2.10]），即得本引理。
\end{proof}

为陈述下面的引理，回忆 [55, 3.1.3]：对 $(-)=\mathcal C$ 或 $\mathbf{Ani}(\mathcal C)$，遗忘函子 $\mathbf{Alg}(-)\to(-)$ 有左伴随 $\mathrm{Fr}_{(-)}$，由自由代数构造给出。

\begin{lem}{2.4}\label{lem:a2-2-4}
对每个 $X\in\mathcal C^{cp}$，$\mathrm{Fr}_{\mathcal C}(X)$ 在函子 $\mathbf{Alg}(\mathcal C)\to\mathbf{Alg}(\mathbf{Ani}(\mathcal C))$ 下的像典范同构于 $\mathrm{Fr}_{\mathbf{Ani}(\mathcal C)}(X)$。
\end{lem}

我们指出，这个引理是 $E_1$-代数特有的，因为 $E_\infty$-代数的类似陈述一般已知不成立。\footnote{我们感谢 P.~Scholze 指出这一点。}

\begin{proof}
我们把 $\mathrm{Fr}_{\mathcal C}(X)$ 视为 $\mathbf{Alg}(\mathbf{Ani}(\mathcal C))$ 中的对象。则由伴随给出典范态射 $\mathrm{Fr}_{\mathbf{Ani}(\mathcal C)}(X)\to\mathrm{Fr}_{\mathcal C}(X)$。为证明它是同构，可以施加遗忘函子 $\mathbf{Alg}(\mathbf{Ani}(\mathcal C))\to\mathbf{Ani}(\mathcal C)$，因为该函子是保守的（[55, 3.2.2.6]）。而在 $\mathbf{Ani}(\mathcal C)$ 中，两个对象都由 $\bigsqcup_{n\ge 0}X^{\otimes n}$ 给出，这只需把 [55, 3.1.3.13] 与「嵌入 $\mathcal C^{cp}\to\mathbf{Ani}(\mathcal C)$ 是幺半的且保持有限余积」这一事实结合起来。
\end{proof}

下面是我们要用的推论。它可以看作 $\mathbf{Mon}(\mathbf{Spc})$ 中对象的一种典范「投射消解」。密切相关的陈述见 [33, 2.1.5]（证明不同）。\ysrc{44}

\begin{cor}{2.5}\label{cor:a2-2-5}
同构 (2.1) 在 $\mathbf{Mon}(\mathbf{Spc})$ 中成立。特别地，对每个 $X_\bullet\in\mathbf{Mon}(\mathbf{Spc})$，
\begin{equation}\tag{2.2}
\begin{aligned}
\Map_{\mathbf{Mon}(\mathbf{Spc})}(\Gamma,X_\bullet)
&=\varprojlim_{(\mathbf{FFM}/\Gamma)^{\mathrm{op}}}\Map_{\mathbf{Mon}(\mathbf{Spc})}(\mathrm{FM}(I),X_\bullet)\\
&=\varprojlim_{(\mathbf{FFM}/\Gamma)^{\mathrm{op}}}X([1])^I .
\end{aligned}
\end{equation}
\end{cor}

当然，(2.1) 对每个 $\Gamma\in\mathbf{Mon}(\mathbf{Spc})$ 都成立，只不过此时 $\mathbf{FFM}/\Gamma$ 可能不再是普通范畴。

\begin{rem}{2.6}\label{rem:a2-2-6}
上述讨论有若干变体，把 $\mathcal C$ 中的幺半群对象换成群对象或半群对象。按照 [55, 5.2.6.2, 4.1.2.12]，我们把群对象视为群状幺半群对象，把半群对象视为非单位幺半群对象，并分别以 $\mathbf{Mon}^{\mathrm{gp}}(\mathcal C)$ 与 $\mathbf{Mon}^{\mathrm{nu}}(\mathcal C)$ 记相应的范畴（当 $\mathcal C=\mathbf{Sets}$ 时省略 $\mathcal C$）。对 $?=\mathrm{gp}$ 或 $\mathrm{nu}$，$\mathbf{Mon}^{?}$ 的紧投射对象仍是有限自由生成的那些。沿用 [76]，我们分别以 $\mathbf{FFG}$ 与 $\mathbf{FFS}$ 记相应的子范畴。我们仍有 $\mathbf{Ani}(\mathbf{Mon}^{?})\cong\mathbf{Mon}^{?}(\mathbf{Spc})$，从而有类似的推论 2.5。事实上，半群情形可类似证明，群情形则由引理 2.3 与 [55, 5.2.6.4] 推出（其实已包含在 [55, 5.2.6.10, 5.2.6.21] 中）。

有自然的遗忘函子
\[
\mathbf{Mon}^{\mathrm{gp}}(\mathbf{Spc})\to\mathbf{Mon}(\mathbf{Spc})\to\mathbf{Mon}^{\mathrm{nu}}(\mathbf{Spc}).
\]
第一个函子与复合函子都是全忠实的。在本文的应用中，我们主要考虑群之间的映射空间，故可在上述三个范畴中的任一个里计算。
\end{rem}

\begin{rem}{2.7}\label{rem:a2-2-7}
存在自然的忠实函子 $\mathbf{Col}:\Delta\to\mathbf{FFM}$，它把 $[n]$ 送到 $\mathrm{FM}(\{x_1,\dots,x_n\})$，把 $f:[n]\to[m]$ 送到由
\[
\mathbf{Col}(f)(x_i)=
\begin{cases}
y_{f(i)}y_{f(i)+1}\cdots y_{f(i+1)-1} & f(i+1)>f(i),\\[0.2em]
e & f(i+1)=f(i)
\end{cases}
\]
定义的 $\mathbf{Col}(f):\mathrm{FM}(\{x_1,\dots,x_n\})\to\mathrm{FM}(\{y_1,\dots,y_m\})$。不难看出 $\mathbf{Col}$ 是共尾的。此外，$\mathbf{Col}^{\mathrm{op}}$ 诱导函子
\[
\mathbf{Fun}(\mathbf{FFM}^{\mathrm{op}},\mathbf{Spc})\to\mathbf{Fun}(\Delta^{\mathrm{op}},\mathbf{Spc}),
\]
它限制为引理 2.3 中的等价。
\end{rem}

\subsection{导出表示空间}

我们固定交换环 $\Lambda$。以 $\mathbf{CAlg}_\Lambda^\heartsuit$ 记交换 $\Lambda$-代数的（普通）范畴，我们有时把 $\mathbf{CAlg}_\Lambda^\heartsuit$ 中的对象称为经典 $\Lambda$-代数。令 $\mathbf{CAlg}_\Lambda=\mathbf{Ani}(\mathbf{CAlg}_\Lambda^\heartsuit)$ 为其导出范畴，并按 Clausen–Scholze 的叫法，称 $\mathbf{CAlg}_\Lambda$ 中的对象为动态化 $\Lambda$-代数（animated $\Lambda$-algebra）。\footnote{这一范畴在 [56, §25] 中记作 $\mathbf{CAlg}_\Lambda^\Delta$，那里它的对象传统上称为单纯 $\Lambda$-代数。不过我们把记号 $\mathbf{CAlg}_\Lambda^\Delta$ 留给 $\mathbf{CAlg}_\Lambda=\mathbf{Ani}(\mathbf{CAlg}_\Lambda^\heartsuit)$ 中的余单纯对象。} 有自然的遗忘函子
\[
\mathbf{CAlg}_\Lambda=\mathbf{Ani}(\mathbf{CAlg}_\Lambda^\heartsuit)\to\mathbf{Ani}(\mathbf{Mod}_\Lambda^\heartsuit)\cong\mathbf{Mod}_\Lambda^{\le 0},
\]
它是保守的，且保持极限与筛状余极限（结合 [56, 25.1.2.2] 与 [55, 3.2.2.1, 3.2.2.6, 3.2.3.1] 可得）。对动态化 $\Lambda$-代数 $A$，我们记 $\pi_i(A)$ 为其底层 $\Lambda$-模的 $(-i)$ 次上同调。称动态化 $\Lambda$-代数 $A$ 为截断的，若它属于某个 $m<\infty$ 的 ${}_{\le m}\mathbf{CAlg}_\Lambda$，这等价于说 $\pi_i(A)=0$ 对 $i>m$。\ysrc{45}

以 $\mathbf{Aff}_\Lambda$（resp.\ $\mathbf{DAff}_\Lambda$）记 $\mathbf{CAlg}_\Lambda^\heartsuit$（resp.\ $\mathbf{CAlg}_\Lambda$）的反范畴。$\mathbf{Aff}_\Lambda$ 中的对象将称为经典仿射 $\Lambda$-概形，或简称为仿射 $\Lambda$-概形；$\mathbf{DAff}_\Lambda$ 中的对象将称为导出仿射 $\Lambda$-概形，或动态化 $\Lambda$-仿射概形。给定 $A\in\mathbf{CAlg}_\Lambda$，$\mathbf{DAff}_\Lambda$ 中相应的对象照常记作 $\Spec A$；给定 $X\in\mathbf{DAff}_\Lambda$，$\mathbf{CAlg}_\Lambda$ 中相应的对象记作 $\Lambda[X]$，称为 $X$ 上的正则函数环。对 $X=\Spec A$，我们以 ${}^{cl}X$ 记底层的经典仿射概形 $\Spec\pi_0(A)$。我们说仿射 $\Lambda$-概形 $\Spec A$ 是（$m$-）截断的，若 $A$ 是（$m$-）截断的。（注意这与「$\Spec A$ 是 $\mathbf{DAff}_\Lambda$ 中的 $m$-截断对象」不同。）

设 $M$ 是 $\Lambda$ 上的经典仿射平坦幺半群概形。它是 $\mathbf{Mon}(\mathbf{Aff}_\Lambda)$ 中的对象。则由 $M$ 定义的函子 $\mathbf{CAlg}_\Lambda^\heartsuit\to\mathbf{Mon}$ 延拓为（保持筛状余极限的）函子
\[
\mathbf{CAlg}_\Lambda=\mathbf{Ani}(\mathbf{CAlg}_\Lambda^\heartsuit)\to\mathbf{Ani}(\mathbf{Mon})\cong\mathbf{Mon}(\mathbf{Spc}),
\]
仍记作 $M$。把定义展开，对 $A\in\mathbf{CAlg}_\Lambda$，$M(A)\in\mathbf{Mon}(\mathbf{Spc})$ 是由
\[
[n]\in\Delta\mapsto\Map_{\mathbf{CAlg}_\Lambda}(\Lambda[M^n],A)\cong\Map_{\mathbf{CAlg}_\Lambda}(\Lambda[M],A)^n
\]
给出的单纯空间。

\begin{defn}{2.8}\label{defn:a2-2-8}
对 $\Gamma\in\mathbf{Mon}(\mathbf{Spc})$，定义
\begin{equation}\tag{2.3}
\mathcal R_{\Gamma,M}:\mathbf{CAlg}_\Lambda\to\mathbf{Spc},\qquad A\mapsto\Map_{\mathbf{Mon}(\mathbf{Spc})}(\Gamma,M(A)).
\end{equation}
\end{defn}

\begin{rem}{2.9}\label{rem:a2-2-9}
我们的定义与 [69, §3.2] 中的定义等价。设
\[
\mathbf{CAlg}_\Lambda^\Delta=\mathbf{Fun}(\Delta,\mathbf{CAlg}_\Lambda)
\]
为 $\mathbf{CAlg}_\Lambda$ 中余单纯对象的范畴。我们有
\begin{equation}\tag{2.4}
\begin{aligned}
\Map(\Gamma,M(A))&=\Map_{\mathbf{Spc}_\Delta}\bigl(\Gamma^\bullet,\Map_{\mathbf{CAlg}_\Lambda}(\Lambda[M^\bullet],A)\bigr)\\
&\cong\Map_{\mathbf{CAlg}_\Lambda^\Delta}\bigl(\Lambda[M^\bullet],C(\Gamma^\bullet,A)\bigr),
\end{aligned}
\end{equation}
其中 $C(\Gamma^\bullet,A)\in\mathbf{CAlg}_\Lambda^\Delta$ 是表示函子
\[
(\mathbf{CAlg}_\Lambda^\Delta)^{\mathrm{op}}\to\mathbf{Spc},\qquad B^\bullet\mapsto\Map_{\mathbf{Spc}_\Delta}\bigl(\Gamma^\bullet,\Map_{\mathbf{CAlg}_\Lambda}(B^\bullet,A)\bigr)
\]
的对象。容易看出 $C(\Gamma^\bullet,A)$ 的第 $n$ 项是由
\begin{equation}\tag{2.5}
C(\Gamma^n,A):=\varprojlim_{\Gamma^n}A=A^{\Gamma^n}
\end{equation}
给出的动态化 $\Lambda$-代数，即从 $\Gamma^n$ 到 $A$ 的映射构成的 $\Lambda$-代数（这一概念在 $\infty$-范畴框架下的定义见 [54, 5.5.2.6]）。

另一方面，若 $M$ 是群概形从而 $M(A)$ 是群状的，则由 [55, 5.2.6.10, 5.2.6.13]，取（单纯空间的）几何实现诱导等价
\begin{equation}\tag{2.6}
\Map_{\mathbf{Mon}(\mathbf{Spc})}(\Gamma,M(A))\to\Map_{\mathbf{Spc}_*}(|\Gamma|,|M(A)|),
\end{equation}
其中 $\mathbf{Spc}_*$ 记带点空间的 $\infty$-范畴（[55, 1.4.2.5]）。因此，我们的定义也与 [35, §5] 中（加框架）表示导出模空间的定义一致。（几何实现 $|\cdot|$ 在 loc.\ cit.\ 中记作 $B(\cdot)$。）
\end{rem}

利用推论 2.5 中 $\Gamma$ 的「消解」，我们立刻得到 $\mathcal R_{\Gamma,M}$ 的如下表现，它特别蕴含 $\mathcal R_{\Gamma,M}$ 作为导出仿射概形的可表示性。\ysrc{46}

\begin{prop}{2.10}\label{prop:a2-2-10}
存在自然同构
\[
\mathcal R_{\Gamma,M}\cong\varprojlim_{(\mathbf{FFM}/\Gamma)^{\mathrm{op}}}M^I,
\]
其中极限在 $\mathbf{DAff}_\Lambda$ 中取。因此，在 $\mathbf{CAlg}_\Lambda$ 中有典范同构
\begin{equation}\tag{2.7}
\Lambda[\mathcal R_{\Gamma,M}]\cong\varinjlim_{\mathbf{FFM}/\Gamma}\Lambda[M^I].
\end{equation}
\end{prop}

如前所述，即使 $\Gamma$ 是离散的，$\mathbf{FFM}/\Gamma$ 也不是滤过范畴。因此，尽管每个 $\Lambda[M^I]$ 只位于上同调次数零，$\Lambda[\mathcal R_{\Gamma,M}]$ 的情形可能并非如此。

\begin{exam}{2.11}\label{exam:a2-2-11}
若 $\Gamma=\mathrm{FM}(I)$，则 $\mathcal R_{\mathrm{FM}(I),M}\cong{}^{cl}\mathcal R_{\mathrm{FM}(I),M}\cong M^I$。这与直观一致：若 $\Gamma$ 是自由的，则没有施加任何关系，此时 ${}^{cl}\mathcal R_{\Gamma,M}$ 不应有非平凡的导出结构。
\end{exam}

\begin{rem}{2.12}\label{rem:a2-2-12}
(1) 上述结果提示如下推广，它对伪表示的讨论有用。考虑
\[
\mathbf{CAlg}_\Lambda^{\mathbf{FFM}}:=\mathbf{Fun}(\mathbf{FFM},\mathbf{CAlg}_\Lambda),
\]
即 [76] 意义下 $\mathbf{FFM}$-代数的范畴。我们把这个范畴中的对象记作 $A^\bullet$。对有限非空集 $I$，以 $A^I$ 记 $\mathrm{FM}(I)$ 在函子 $A^\bullet$ 下的像。例如，联系于 $M\in\mathbf{Mon}(\mathbf{Aff}_\Lambda)$ 有 $\mathbf{FFM}$-代数 $\Lambda[M]^\bullet$，它把 $\mathrm{FM}(I)$ 送到 $\Lambda[M^I]$。另一方面，设 $B\in\mathbf{CAlg}_\Lambda$，$\Gamma\in\mathbf{Mon}(\mathbf{Spc})$。我们把 $B$ 视为取值 $B$ 的常函子 $\mathbf{FFM}/\Gamma\to\mathbf{CAlg}_\Lambda$。它沿 $\mathbf{FFM}/\Gamma\to\mathbf{FFM}$ 的右 Kan 延拓不是别的，正是 $\mathbf{FFM}$-代数 $\mathrm{FM}(I)\mapsto C(\Gamma^I,B)$。

现在，对 $\mathbf{FFM}$-代数 $A^\bullet$ 与 $\Gamma\in\mathbf{Mon}(\mathbf{Spc})$，我们可以定义
\[
\mathcal R_{\Gamma,\Spec A^\bullet}:=\varprojlim_{\mathrm{FM}(I)\in(\mathbf{FFM}/\Gamma)^{\mathrm{op}}}\Spec A^I,
\quad\text{于是}\quad
\Lambda[\mathcal R_{\Gamma,\Spec A^\bullet}]=\varinjlim_{\mathbf{FFM}/\Gamma}A^I .
\]
当 $A^\bullet=\Lambda[M]^\bullet$ 时，$\mathcal R_{\Gamma,\Spec A^\bullet}$ 还原 $\mathcal R_{\Gamma,M}$。

注意对每个 $B\in\mathbf{CAlg}_\Lambda$，有
\begin{equation}\tag{2.8}
\begin{aligned}
\Map_{\mathbf{CAlg}_\Lambda}\bigl(\Lambda[\mathcal R_{\Gamma,\Spec A^\bullet}],B\bigr)
&\cong\Map_{\mathbf{CAlg}_\Lambda^{\mathbf{FFM}/\Gamma}}(A^\bullet,B)\\
&\cong\Map_{\mathbf{CAlg}_\Lambda^{\mathbf{FFM}}}\bigl(A^\bullet,C(\Gamma^\bullet,B)\bigr).
\end{aligned}
\end{equation}
作为推论，我们看到注 2.7 中的函子 $\mathbf{Col}$ 诱导同构
\[
\Map_{\mathbf{CAlg}_\Lambda^{\mathbf{FFM}}}\bigl(\Lambda[M]^\bullet,C(\Gamma^\bullet,B)\bigr)
\to\Map_{\mathbf{CAlg}_\Lambda^\Delta}\bigl(\Lambda[M]^\bullet,C(\Gamma^\bullet,B)\bigr).
\]
我们不知道，若把 $\Lambda[M]^\bullet$ 与 $C(\Gamma^\bullet,B)$ 换成更一般的 $\mathbf{FFM}$-代数，这是否仍成立。

(2) 可以把 $\mathbf{FFM}$ 换成 [76] 中考虑的 $\mathbf{FFS}$ 或 $\mathbf{FFG}$。我们不再重复此类注记。
\end{rem}

让我们回到 $\mathcal R_{\Gamma,M}$，讨论其上的一些向量丛。为简单起见，从现在起我们假设 $\Gamma$ 是离散的，即 $\mathbf{Mon}$ 中的对象。这对我们的目的已经足够，并简化下面的讨论。如前述，我们把它等同于只有一个对象的范畴，再经由 Milnor 构造等同于一个单纯集。\ysrc{47}

关于动态化环上的模论，我们引用 [56, §25.2.1]（进一步的阐述见 [14, 5.1]）。对动态化 $\Lambda$-代数 $A$，以 $\mathbf{Mod}_A$ 记 $A$-模的 $\infty$-范畴，$\mathbf{Mod}_A^{\le 0}$ 记连通对象构成的全子范畴。若 $A$ 是经典的，$\mathbf{Mod}_A^{\le 0}$ 也等价于前面引入的 $\mathbf{Ani}(\mathbf{Mod}_A^\heartsuit)$。我们也把 $A$-模称为 $\Spec A$ 上的拟凝聚层。

现在，对 $M$ 在有限投射 $\Lambda$-模上的表示 $W$，以 ${}_\Gamma W$ 记 $\mathcal R_{\Gamma,M}$ 上的（平凡）向量丛 $\Lambda[\mathcal R_{\Gamma,M}]\otimes_\Lambda W$。我们将给 ${}_\Gamma W$ 装备 $\Gamma$ 的作用，更确切地说，构造 $\mathbf{Mon}(\mathbf{Spc})$ 中的典范态射
\begin{equation}\tag{2.9}
\Gamma\to\End({}_\Gamma W).
\end{equation}
这里 $\End({}_\Gamma W)\in\mathbf{Mon}(\mathbf{Spc})$ 记 ${}_\Gamma W$（视为 $\mathcal R_{\Gamma,M}$ 上的连通拟凝聚层）的导出自同态环（例 2.2）。

下文中，为简单起见我们把 ${}_{\mathrm{FM}(I)}W$ 记作 ${}_I W$。注意在 $\mathbf{Mon}(\mathbf{Spc})$ 中有典范同构 $\varinjlim_{\mathbf{FFM}/\Gamma}\End({}_I W)\to\End({}_\Gamma W)$。于是由推论 2.5，只需对每个 $u:\mathrm{FM}(I)\to\Gamma$ 构造态射 $\mathrm{FM}(I)\to\End({}_I W)$，并与 $\mathbf{FFM}/\Gamma$ 中的态射相容。我们指出，后一相容性可在普通范畴层面验证。

其次，经由包含 $\{i\}\subset I$，只需假设 $I=\{1\}$，并在 $M$ 上构造 ${}_{\{1\}}W$ 的自同态，即 $\Lambda[M]\otimes W$ 的 $\Lambda[M]$-线性自同态。但这正是余作用映射
\begin{equation}\tag{2.10}
\mathrm{coact}:W\to\Lambda[M]\otimes_\Lambda W .
\end{equation}
这就完成了 (2.9) 的构造。

\begin{rem}{2.13}\label{rem:a2-2-13}
(1) 这里对 $\Gamma$ 在 ${}_\Gamma W$ 的纤维上的作用 (2.9) 给出更具体的描述。表示 $W$ 诱导 $\Lambda$ 上幺半群概形的同态 $M\to\End(W)$，其中 $\End(W)(A)=\End_{\mathbf{Mod}_A^{\le 0}}(W\otimes A)\in\mathbf{Mon}(\mathbf{Spc})$。设 $\Spec A\to\mathcal R_{\Gamma,M}$ 是 $\mathcal R_{\Gamma,M}$ 的一个点，对应于同态 $\rho:\Gamma\to M(A)$。${}_\Gamma W$ 在 $\rho$ 上的纤维通常记作 $W_\rho$，它就是 $W\otimes_\Lambda A$，$\Gamma$ 经由 $\Gamma\xrightarrow{\rho}M(A)\to\End(W)(A)$ 作用其上。在 (2.4) 中，我们把 $\rho$ 解释为余单纯代数的映射 $\Lambda[M^\bullet]\to C(\Gamma^\bullet,A)$。本着同样的精神，我们也可以把这个作用解释为 $C(\Gamma^\bullet,A)$（从而 $A$）上的余单纯模 $C(\Gamma^\bullet,W_\rho)$，方式如下。余作用 (2.10) 延拓为 $\Lambda[M^\bullet]$ 上的余单纯模 $\Lambda[M^\bullet]\otimes_\Lambda W$。而 $C(\Gamma^\bullet,W_\rho)$ 是它沿 $\rho$ 的基变换。

(2) 若 $W$ 是 $M^J$ 的表示（$J$ 为有限集），则 ${}_\Gamma W$ 带有 $\Gamma^J$ 的作用：先对 $\mathcal R_{\Gamma^J,M^J}$ 施加上述构造，再沿态射 $\mathcal R_{\Gamma,M}\to\mathcal R_{\Gamma^J,M^J}$ 把 ${}_{\Gamma^J}W$ 上的 $\Gamma^J$-作用拉回。
\end{rem}

我们可以把 (2.9) 解释为从 $\Gamma$ 到 $\mathcal R_{\Gamma,M}$ 上拟凝聚层范畴的函子，它把 $\Gamma$ 的唯一对象送到 ${}_\Gamma W$（见例 2.2）。

\begin{defn}{2.14}\label{defn:a2-2-14}
$\Gamma$ 以 $W$ 为系数的「泛」同调是 $\mathcal R_{\Gamma,M}$ 上由
\[
C_*(\Gamma,{}_\Gamma W):=\varinjlim_\Gamma{}_\Gamma W
\]
定义的拟凝聚层复形。
\end{defn}

由于张量积保持余极限，$C_*(\Gamma,{}_\Gamma W)$ 沿注 2.13 中由 $\rho:\Gamma\to M(A)$ 给出的 $\Spec A\to\mathcal R_{\Gamma,M}$ 的（导出）拉回，正是 $\mathbf{Mod}_A^{\le 0}$ 中计算 $\varinjlim_\Gamma W_\rho$ 的复形。若 $A$ 是经典的，这不过是 $\Gamma$ 以 $W_\rho$ 为系数的通常同调。\ysrc{48}

有典范同构
\begin{equation}\tag{2.11}
C_*(\Gamma,{}_\Gamma W)\cong\varinjlim_{\mathbf{FFM}/\Gamma}\Lambda[\mathcal R_{\Gamma,M}]\otimes_{\Lambda[M^I]}C_*(\mathrm{FM}(I),{}_I W),
\end{equation}
它由推论 2.5 构造，
\begin{align*}
\varinjlim_\Gamma{}_\Gamma W&\cong\varinjlim_{\mathbf{FFM}/\Gamma}\varinjlim_{\mathrm{FM}(I)}\Lambda[\mathcal R_{\Gamma,M}]\otimes_{\Lambda[M^I]}{}_I W\\
&\cong\varinjlim_{\mathbf{FFM}/\Gamma}\Lambda[\mathcal R_{\Gamma,M}]\otimes_{\Lambda[M^I]}\varinjlim_{\mathrm{FM}(I)}{}_I W .
\end{align*}

考虑 $C_*$ 的约化版本是方便的。按定义，有自然映射 ${}_\Gamma W\to C_*(\Gamma,{}_\Gamma W)$。我们把它在 $\mathcal R_{\Gamma,M}$ 上拟凝聚层范畴中的纤维记作 $\overline{C}_*(\Gamma,{}_\Gamma W)[-1]$，于是有特异三角
\begin{equation}\tag{2.12}
\overline{C}_*(\Gamma,{}_\Gamma W)[-1]\to{}_\Gamma W\to C_*(\Gamma,{}_\Gamma W)\to .
\end{equation}
此时 (2.11) 在把 $C_*$ 换成 $\overline{C}_*$ 后仍成立。考虑约化版本的好处在于我们有如下典范同构
\begin{equation}\tag{2.13}
{}_I W^{\oplus I}\cong\overline{C}_*(\mathrm{FM}(I),{}_I W)[-1],
\end{equation}
它由自由幺半群同调的下述两项复形（位于上同调次数 $[-1,0]$）的计算得到：
\[
\bigoplus_{i\in I}{}_I W\xrightarrow{\ \oplus_{i\in I}(\gamma_i-1)\ }{}_I W,
\]
其中 $\gamma_i$ 记 $\mathrm{FM}(I)$ 中对应于 $i\in I$ 的生成元。特别地，$\overline{C}_*(\mathrm{FM}(I),{}_I W)[-1]$ 位于 $\mathcal R_{\mathrm{FM}(I),M}\cong M^I$ 上拟凝聚层的阿贝尔范畴中。

现设 $f:\mathrm{FM}(I)\to\mathrm{FM}(J)$ 是幺半群态射。它诱导同调之间的态射
\[
\Lambda[M^J]\otimes_{\Lambda[M^I]}\overline{C}_*(\mathrm{FM}(I),{}_I W)[-1]
\to\overline{C}_*(\mathrm{FM}(J),{}_J W)[-1].
\]
在同构 (2.13) 之下，它由 $\Lambda[M^I]$-线性映射
\begin{equation}\tag{2.14}
{}_I W^{\oplus I}\to{}_J W^{\oplus J}
\end{equation}
给出，我们现在更明确地描述它。注意每个这样的 $f:\mathrm{FM}(I)\to\mathrm{FM}(J)$ 都是以下两类映射的复合：
\begin{itemize}[leftmargin=1.6em,itemsep=0.2em]
\item $f$ 把 $\mathrm{FM}(I)$ 的生成元送到 $\mathrm{FM}(J)$ 的生成元或单位元，即 $f$ 由带点集之间的映射 $I\cup\{*\}\to J\cup\{*\}$ 诱导；
\item $f:\mathrm{FM}(\{1,\dots,n\})\to\mathrm{FM}(\{1,\dots,n+1\})$ 把 $\gamma_i\mapsto\gamma_i$（$i\le n-1$），且 $f(\gamma_n)=\gamma_n\gamma_{n+1}$。
\end{itemize}

因此，只需分别在这两种情形下理解 (2.14)。把 (2.13) 的构造展开，我们看到第一种情形由
\begin{equation}\tag{2.15}
\begin{aligned}
(w_i)_{i\in I}\in{}_I W^{\oplus I}&\mapsto(v_j)_{j\in J}\in{}_J W^{\oplus J},\\
v_j&=\sum_{i\in f^{-1}(j)}1\otimes w_i,
\end{aligned}
\end{equation}
给出，而第二种情形由
\begin{equation}\tag{2.16}
\begin{aligned}
(w_i)\in{}_{\{1,\dots,n\}}W^{\oplus n}&\mapsto(v_j)\in{}_{\{1,\dots,n+1\}}W^{\oplus(n+1)},\\
v_i&=1\otimes w_i,\ i\le n,\qquad v_{n+1}=\gamma_n(1\otimes w_n).
\end{aligned}
\end{equation}
给出。\ysrc{49}

现在，当 $M$ 是 $\Lambda$ 上的仿射光滑群概形时，我们可以计算 $\mathcal R_{\Gamma,M}$ 上的余切复形。以 $\mathrm{Ad}^*$ 记 $M$ 在其李代数 $\mathfrak m$ 的对偶上的余伴随表示。

我们回忆，对动态化 $\Lambda$-代数 $A$，（代数）余切复形 $\mathbb L_A$ 是连通的 $A$-模，使得对每个 $A\to B$ 与连通的 $B$-模 $V$，
\[
\Map_{\mathbf{Mod}_A^{\le 0}}(\mathbb L_A,V)\cong\Map_{\mathbf{CAlg}_{\Lambda/B}}(A,B\oplus V),
\]
其中 $B\oplus V\to B$ 记 $B$ 被 $V$ 在 $\mathbf{CAlg}_\Lambda$ 中的平凡平方零扩张，$\mathbf{CAlg}_{\Lambda/B}$ 记带有到 $B$ 的 $\Lambda$-代数映射的动态化 $\Lambda$-代数的范畴。详见 [56, 25.3.1, 25.3.2]。若 $A$ 是经典光滑 $\Lambda$-代数，则 $\mathbb L_A\cong\pi_0(\mathbb L_A)=\Omega_A$ 就是 $A$ 的 Kähler 微分。若 $A\to B$ 是 $\mathbf{CAlg}_\Lambda$ 中的态射，则在 $\mathbf{Mod}_B^{\le 0}$ 中有自然态射 $B\otimes_A\mathbb L_A\to\mathbb L_B$，相对余切复形 $\mathbb L_{B/A}$ 定义为它的纤维。

\begin{prop}{2.15}\label{prop:a2-2-15}
假设 $M$ 是 $\Lambda$ 上的仿射光滑群概形。对每个 $\Gamma$，$\mathcal R_{\Gamma,M}$ 的余切复形典范同构于 $\overline{C}_*(\Gamma,{}_\Gamma\mathrm{Ad}^*)[-1]$。
\end{prop}

\begin{proof}
注意若 $A=\varinjlim A_i$ 是 $\mathbf{CAlg}_\Lambda$ 中的余极限，则
\begin{equation}\tag{2.17}
\mathbb L_A\cong\varinjlim\,(A\otimes_{A_i}\mathbb L_{A_i}).
\end{equation}
把它应用于 $\Lambda[\mathcal R_{\Gamma,M}]=\varinjlim_{\mathbf{FFM}/\Gamma}\Lambda[M^I]$。把 (2.11) 与 (2.17) 比较可知，只需对每个 $f:\mathrm{FM}(I)\to\mathrm{FM}(J)$ 建立下图（在 $\Lambda[M^J]$-模的阿贝尔范畴中）的交换性：
\begin{equation}\tag{2.18}
\begin{tikzcd}[column sep=small,row sep=small]
\Lambda[M^J]\otimes_{\Lambda[M^I]}({}_I\mathrm{Ad}^*)^{\oplus I}\arrow[r]\arrow[d,"\cong"'] & ({}_J\mathrm{Ad}^*)^{\oplus J}\arrow[d,"\cong"]\\
\Lambda[M^J]\otimes_{\Lambda[M^I]}\Omega_{M^I/\Lambda}\arrow[r] & \Omega_{M^J}.
\end{tikzcd}
\end{equation}
现在若把 $\Omega_M$ 与 $\Lambda[M]\otimes\mathrm{Ad}^*$ 等同起来（把 $\mathrm{Ad}^*$ 视为左不变微分的空间），则竖直同构变得清楚，而图的交换性由 (2.15) 与 (2.16) 中展示的显式计算得到。
\end{proof}

\begin{rem}{2.16}\label{rem:a2-2-16}
有时取 $\mathcal R_{\Gamma,M}$ 的余切复形的线性对偶是方便的。给定 $\rho:\Gamma\to M(A)$，$\mathcal R_{\Gamma,M}$ 在 $\rho$ 处的切空间 $T_\rho\mathcal R_{\Gamma,M}$ 是 $\mathbb L_{\mathcal R_{\Gamma,M}}|_\rho$（视为 $\mathbf{Mod}_A$ 中的对象）的 $A$-线性对偶，它同构于 $\overline{C}^*(\Gamma,\mathrm{Ad}_\rho)[1]$。这里
\[
C^*(\Gamma,\mathrm{Ad}_\rho):=\varprojlim_{\Gamma}\mathrm{Ad}_\rho,
\]
其中极限在 $\mathbf{Mod}_A$ 中取，而 $\overline{C}^*(\Gamma,\mathrm{Ad}_\rho)[1]$ 是其约化版本，即 $C^*(\Gamma,\mathrm{Ad}_\rho)\to\mathrm{Ad}_\rho$ 的余纤维。若 $A$ 是经典的，这就是 $\Gamma$ 以 $M$ 的伴随表示 $\mathrm{Ad}$ 为系数的通常上同调。注意对 $M$ 的表示 $W$，$C^*(\Gamma,W_\rho)$ 可以等同于注 2.13(1) 中余单纯 $A$-模 $C(\Gamma^\bullet,W_\rho)$ 的全化。\ysrc{50}
\end{rem}

让我们转到下一个话题。注意若 $\Gamma$ 有限生成且 $\Lambda$ 诺特，则非导出空间 ${}^{cl}\mathcal R_{\Gamma,M}$ 在 $\Lambda$ 上是有限型的。事实上，选取满射 $\mathrm{FM}(I)\to\Gamma$，${}^{cl}\mathcal R_{\Gamma,M}$ 就实现为 ${}^{cl}\mathcal R_{\mathrm{FM}(I),M}\cong M^I$ 的闭子概形。现在我们讨论 $\mathcal R_{\Gamma,M}$ 的类似陈述。

回忆对紧生成 $\infty$-范畴 $\mathcal C$，对象 $c$ 称为几乎紧的，若对每个 $n\ge 0$，$\tau_{\le n}c$ 在 ${}_{\le n}\mathcal C$ 中是紧的（[55, 7.2.4.8]）。$\mathbf{CAlg}_\Lambda$ 中的几乎紧对象也称为几乎有限表现的；对动态化 $\Lambda$-代数 $A$，$\mathbf{Mod}_A^{\le 0}$ 中的几乎紧对象也称为连通几乎完美 $A$-模。若 $\Lambda$ 诺特，则 $A$ 在 $\Lambda$ 上几乎有限表现当且仅当 $\pi_0(A)$ 是有限生成 $\Lambda$-代数且每个 $\pi_i(A)$ 是有限生成 $\pi_0(A)$-模（[57, 3.1.5]）。特别地，若 $\Lambda$ 诺特，则有限型的经典 $\Lambda$-代数在视为动态化 $\Lambda$-代数时是几乎有限表现的。

另一方面，回忆称群（甚至幺半群）$\Gamma$ 是 $FP_\infty(k)$ 型的，若平凡 $k\Gamma$-模容许一个消解 $P^\bullet\to k$，其每一项都是有限投射 $k\Gamma$-模，其中 $k\Gamma$ 记 $\Gamma$ 的群（或幺半群）代数。例如，有限群总是 $FP_\infty(k)$ 型的。更一般地，若 $\Gamma$ 的分类空间可实现为在每个次数 $n\ge 0$ 上只有有限个胞腔的 CW 复形（这样的群称为 $F_\infty$ 型的），则 $\Gamma$ 是 $FP_\infty(k)$ 型的。

\begin{prop}{2.17}\label{prop:a2-2-17}
假设 $\Lambda$ 诺特，$M$ 是 $\Lambda$ 上的光滑仿射群概形。若 $\Gamma$ 有限生成且为 $FP_\infty(k)$ 型，则 $\mathcal R_{\Gamma,M}$ 在 $\Lambda$ 上几乎有限表现。
\end{prop}

\begin{proof}
由于 $\Gamma$ 有限生成，${}^{cl}\mathcal R_{\Gamma,M}$ 是有限型的。利用 [57, 3.2.18] 与命题 2.15，只需证明 $\overline{C}_*(\Gamma,{}_\Gamma\mathrm{Ad}^*)[-1]$ 是几乎完美的。由于 $\Gamma$ 是 $FP_\infty(k)$ 型的，这个复形到每个经典 $\Lambda$-代数 $A$ 上的拉回是每一项都为有限投射 $A$-模的连通复形，从而是几乎完美的。由 [56, 2.7.3.2]，这意味着 $\overline{C}_*(\Gamma,{}_\Gamma\mathrm{Ad}^*)[-1]$ 是几乎完美的。
\end{proof}

\begin{rem}{2.18}\label{rem:a2-2-18}
还有更精细的概念，例如 $n$ 阶有限生成的动态化 $\Lambda$-代数与 $FP_n(k)$ 型的群。可以用这些概念表述上述命题的一个精细版本。
\end{rem}

\begin{prop}{2.19}\label{prop:a2-2-19}
假设同命题 2.17。以 $d$ 记 $M$ 在 $\Lambda$ 上的相对维数。此外，假设对每个由表示 $\rho:\Gamma\to M(\kappa)$ 给出的域值点 $\Spec\kappa\to\mathcal R_{\Gamma,M}$，都有
\[
H_i(\Gamma,\mathrm{Ad}^*_\rho)=0\ \text{对}\ i>2,\qquad\text{且}\qquad
\dim_\kappa{}^{cl}\mathcal R_{\Gamma,M}\le d-\dim(-1)^iH_i(\Gamma,\mathrm{Ad}^*_\rho),
\]
其中\tn{上一显示式末项作 $d-\dim(-1)^iH_i(\Gamma,\mathrm{Ad}^*_\rho)$，而命题 2.19 的证明中给出的等式是 $d-\sum_i(-1)^i\dim H_i(\Gamma,\mathrm{Ad}^*_\rho)$，二者不一致，疑为底本排版时漏去 $\sum_i$、并使 $\dim$ 的位置错乱所致。译文照录原文。} $\dim_\kappa{}^{cl}\mathcal R_{\Gamma,M}$ 记 ${}^{cl}\mathcal R_{\Gamma,M}$ 在 $\kappa$ 处相对于 $\Lambda$ 的维数。则 $\mathcal R_{\Gamma,M}={}^{cl}\mathcal R_{\Gamma,M}$ 是局部完全交。此时，$\mathcal R_{\Gamma,M}$ 在几何点 $\rho\in\mathcal R_{\Gamma,M}$ 处光滑，当且仅当 $\mathcal R_{\Gamma,M}$ 在 $\rho$ 处关于 $\Lambda$ 平坦且 $H_2(\Gamma,\mathrm{Ad}^*_\rho)=0$。
\end{prop}

\begin{proof}
按我们的假设，$\mathcal R_{\Gamma,M}$ 在 $\Lambda$ 上几乎有限表现，且其余切复形的 Tor-幅度 $\le 1$。所以它在 [57, 3.4.15] 的意义下是拟光滑的（当 $\Lambda$ 是特征零的域时也见 [2, 2.1.3]）。选取满射 $\mathrm{FM}(I)\to\Gamma$，诱导态射 $\mathcal R_{\Gamma,M}\to\mathcal R_{\mathrm{FM}(I),M}$。由 loc.\ cit.\ 中的论证可知，在 $M^I$ 上 Zariski 局部地——即把 $M^I$ 换成一个开子概形 $\Spec A\subset M^I$、把 $\mathcal R_{\Gamma,M}$ 换成 $\Spec B:=\Spec A\times_{M^I}\mathcal R_{\Gamma,M}$ 之后——存在态射 $\Spec A\to\mathbb A^m:=\Spec\Lambda[x_1,\dots,x_m]$，使得 $\Spec B\cong\Spec A\times_{\mathbb A^m}\{0\}$。特别地，在 $\Spec B$ 的每个域值点 $\kappa$ 处有 $\dim_\kappa{}^{cl}\mathcal R_{\Gamma,M}\ge\dim_\kappa M^I-m$。另一方面，特异三角 $B\otimes_A\mathbb L_A\to\mathbb L_B\to\mathbb L_{B/A}$ 蕴含对 $\Spec B$ 的每个点 $\kappa$ 有
\[
\dim_\kappa M^I-m=d-\sum_i(-1)^i\dim H_i(\Gamma,\mathrm{Ad}^*_\rho).
\]
由我们的假设推出 $\dim_\kappa{}^{cl}\mathcal R_{\Gamma,M}=\dim_\kappa M^I-m$。这意味着 $\mathcal R_{\Gamma,M}={}^{cl}\mathcal R_{\Gamma,M}$ 是局部完全交。

最后，$\mathcal R_{\Gamma,M}$ 在 $\rho$ 处光滑当且仅当它平坦且 $\dim(\Omega_{\mathcal R_{\Gamma,M}}\otimes\kappa)=\dim_\kappa\mathcal R_{\Gamma,M}$。但由上面的等式，后一条件等价于 $H_2(\Gamma,\mathrm{Ad}^*_\rho)=0$。
\end{proof}

到目前为止，我们关注的是所谓的加框架表示空间。我们也简要讨论表示叠。首先，所谓 $\Lambda$ 上的预叠，指的是一个（可及的）\footnote{这是一个集合论假设（见 [54, 5.4.2.5]）。或者，我们可以对所受考虑的代数的大小加以限制。}函子 $\mathcal F:\mathbf{CAlg}_\Lambda\to\mathbf{Spc}$。$\Lambda$ 上所有预叠构成 $\infty$-范畴 $\mathbf{Fun}(\mathbf{CAlg}_\Lambda,\mathbf{Spc})$。称预叠为叠，若它关于 $\mathbf{CAlg}_\Lambda$ 上的平展拓扑是层。我们以 $\mathbf{Shv}(\mathbf{CAlg}_\Lambda)$ 记 $\mathbf{Fun}(\mathbf{CAlg}_\Lambda,\mathbf{Spc})$ 中由叠构成的全子范畴。与经典情形一样，经由 Yoneda 嵌入，$\mathbf{DAff}_\Lambda$ 构成 $\mathbf{Shv}(\mathbf{CAlg}_\Lambda)$ 的全子范畴。$\Lambda$ 上的导出 Artin 叠是满足某些性质的叠。对（预）叠 $\mathcal F$，以 ${}^{cl}\mathcal F$ 记它在经典 $\Lambda$-代数上的限制，称为其底层的经典（预）叠。注意若 $\mathcal F=\Spec A$，则 ${}^{cl}\mathcal F$ 由 $\Spec\pi_0(A)$ 表示，这与我们之前对 ${}^{cl}\Spec A$ 的定义一致。精确的定义与进一步讨论见 [57, §5]。

现假设存在 $\Lambda$ 上的光滑仿射群概形 $H$，它通过幺半群自同构作用在 $M$ 上。这给出 $\mathbf{Mon}(\mathbf{Aff}_\Lambda)$ 中的单纯对象：把 $[n]\in\Delta$ 送到 $H^n\times M$（幺半结构来自 $M$），并按通常方式给出由作用映射与投影映射得来的各种面映射。于是施加构造 (2.3) 得到单纯导出仿射概形（省略退化映射）
\begin{equation}\tag{2.19}
\cdots\rightrightarrows H\times H\times\mathcal R_{\Gamma,M}\rightrightarrows H\times\mathcal R_{\Gamma,M}\rightrightarrows\mathcal R_{\Gamma,M},
\end{equation}
这相当于 $H$ 在 $\mathcal R_{\Gamma,M}$ 上的一个作用。\ysrc{51}

\begin{defn}{2.20}\label{defn:a2-2-20}
令 $\mathcal R_{\Gamma,M\sslash H}:=\mathcal R_{\Gamma,M}/H$ 为上述 $H$-作用的商叠，即 (2.19) 在 $\mathbf{Shv}(\mathbf{CAlg}_\Lambda)$ 中的几何实现。若 $M=H$ 且 $H$ 通过共轭作用，我们记 $\mathcal X_{\Gamma,H}$ 为 $\mathcal R_{\Gamma,H\sslash H}$，并称之为 $\Gamma$ 的 $H$-表示叠。
\end{defn}

\begin{rem}{2.21}\label{rem:a2-2-21}
显然 ${}^{cl}\mathcal X_{\Gamma,H}$ 是文献中研究的通常表示叠。特别地，对代数闭域 $\kappa$，$\mathcal X_{\Gamma,H}$ 的 $\kappa$-点分类
\end{rem}

% 说明：注 2.21 的末句在底本第 51 页（PDF 64）页末被截断（原文止于「the κ-points of X_{Γ,H} classify」），
% 续文见下一 chunk（ch-a2-02.tex）开头。此处已闭合 rem 环境，以免环境跨文件悬挂。
